\section{Độ tin cậy, kiểm thử và trường hợp biên}

\subsection{Tích hợp hệ thống đầu-cuối và kiểm thử}

\subsubsection{Tích hợp hệ thống đầu-cuối và kiểm thử kỹ lưỡng}

\hspace{0.5cm} Sau khi hoàn thiện các mô hình động học và giải thuật khắc phục khiếm khuyết quỹ đạo, hệ thống bước vào giai đoạn tích hợp đầu-cuối (End-to-End System Integration). Mục tiêu cốt lõi là kết nối toàn bộ các phân tầng kỹ thuật: từ khối trích xuất dữ liệu thực nghiệm, khối sinh nhiễu cảm biến, bộ hợp nhất cảm biến hỗn hợp, chuỗi lọc ước lượng Kalman và Alpha-Beta, kênh truyền thông tin thời gian thực đến giao diện hiển thị bản đồ số và bảng giám sát tài nguyên phần cứng.

Công tác tích hợp bảo đảm dữ liệu đo trắc địa được xử lý tuần tự, liên tục ở tần số 10 Hz với độ trễ thấp ($< 60$ $\mu$s mỗi chu kỳ), hoàn toàn không xảy ra hiện tượng mất mát gói tin, xung đột tài nguyên hay tích lũy sai số trôi dạt.

\subsubsection{Chiến lược kiểm định tích hợp và kiểm thử phân tầng}

\hspace{0.5cm} Dựa trên kiến trúc hệ thống 4 tầng, quy trình quản lý phiên mô phỏng và luồng truyền thông hai chiều WebSocket đã thiết lập tại Chương 2, công tác kiểm thử tích hợp đầu-cuối tập trung đánh giá tính ổn định, độ trễ truyền thông và tính toàn vẹn của dữ liệu bám bắt trên toàn bộ chuỗi mắt xích phần mềm.


\subsubsection{Chiến lược kiểm thử phân tầng theo mô hình kim tự tháp (Test Pyramid)}

\hspace{0.5cm} Toàn bộ hệ thống được bảo vệ bởi chiến lược kiểm thử phân tầng ba lớp nghiêm ngặt, tuân thủ nguyên tắc kim tự tháp kiểm thử phần mềm:

\begin{table}[H]
\centering
\caption{Phân tầng cấu trúc 163 ca kiểm thử tự động toàn diện của hệ thống}
\label{tab:test-pyramid-distribution}
\begin{tabularx}{\textwidth}{>{\RaggedRight\arraybackslash}p{4.2cm} >{\RaggedRight\arraybackslash}X cc}
\toprule
\textbf{Tầng kiểm thử} & \textbf{Phạm vi xác minh kỹ thuật} & \textbf{Số lượng ca} & \textbf{Tỷ trọng (\%)} \\
\midrule
Tầng 1: Đơn vị (Unit Tests) & Giải thuật Kalman, Alpha-Beta, Geodetics, RSS & 98 & 60{,}1\% \\
\addlinespace[3pt]
Tầng 2: Tích hợp (Integration) & Bộ tải dữ liệu, Chuỗi nối tầng LLA-ENU, Ranh giới & 42 & 25{,}8\% \\
\addlinespace[3pt]
Tầng 3: Đầu-cuối (End-to-End) & Luồng REST-WS, Xuất tệp CSV, Dashboard Telemetry & 23 & 14{,}1\% \\
\midrule
\textbf{Tổng cộng toàn hệ thống} & \textbf{Toàn diện 3 tầng kiến trúc} & \textbf{163} & \textbf{100{,}0\% (Đạt 100\%)} \\
\bottomrule
\end{tabularx}
\end{table}

Chi tiết về cấu trúc 163 ca kiểm thử tự động, hệ thống được phân bổ như sau: 125 ca kiểm thử bao phủ toàn diện các mô-đun: bộ lọc Kalman, bộ lọc Alpha-Beta, phép biến đổi tọa độ (Geodetics), hợp nhất cảm biến và cơ chế phản xạ biên mềm. Đồng thời, 38 ca kiểm thử chuyên sâu bảo vệ logic của các trình nạp dữ liệu định tuyến (Geolife, mạng đường giao thông thực tế và động học Drone). Toàn bộ ca kiểm thử liên quan đến nhiễu ngẫu nhiên đều được cố định giá trị khởi tạo (Seed) nhằm duy trì kết quả kiểm thử mang tính xác định (Deterministic) và lặp lại chính xác (Reproducible).


\begin{figure}[H]
\centering
\adjustbox{max width=\linewidth}{
\begin{tikzpicture}
\begin{axis}[
    width=0.85\textwidth, height=7.0cm,
    ybar, bar width=14pt,
    xlabel={Phân hệ chức năng}, ylabel={Độ bao phủ (\%)},
    ymin=0, ymax=100,
    xtick=data,
    xticklabels={Kalman, Alpha-Beta, Geodetics, Fusion, Boundary, Noise, Router, Dashboard},
    xticklabel style={rotate=20, anchor=east, font=\footnotesize},
    ytick={0, 25, 50, 75, 100},
    grid=major, grid style={dashed, gray!30},
    nodes near coords, nodes near coords style={font=\scriptsize},
    legend style={at={(0.5,-0.28)}, anchor=north, legend columns=2, font=\footnotesize},
]
\addplot[fill=blue!50] coordinates {(1,96) (2,94) (3,98) (4,91) (5,92) (6,88) (7,82) (8,85)};
\addplot[fill=green!50] coordinates {(1,93) (2,90) (3,95) (4,87) (5,89) (6,84) (7,78) (8,80)};
\legend{Bao phủ dòng lệnh (Line Coverage), Bao phủ nhánh rẽ (Branch Coverage)}
\end{axis}
\end{tikzpicture}
}
\caption{Mức độ bao phủ kiểm thử tự động theo từng phân hệ chức năng cốt lõi}
\label{fig:coverage-barchart-w11}
\end{figure}

Mức độ bao phủ dòng lệnh trung bình của toàn hệ thống đạt $90{,}8\%$ và bao phủ nhánh rẽ đạt $87{,}0\%$, trong đó các module tính toán toán học cốt lõi (Geodetics, Kalman) đạt trên $95\%$.

\subsubsection{Kiểm tra tải đồng thời và độ ổn định lâu dài (Load and Stress Testing)}

\hspace{0.5cm} Quá trình kiểm thử tích hợp (Integration Testing) không chỉ dừng lại ở việc xác thực các giá trị đầu ra tại một thời điểm, mà còn phải đánh giá độ ổn định của toàn bộ chuỗi xử lý qua hàng nghìn chu kỳ liên tiếp. Chẳng hạn, trong kịch bản bám bắt xe máy di chuyển trên mạng lưới giao thông thực tế, bộ lọc Alpha-Beta phải thể hiện khả năng giảm chấn tối ưu, không tạo ra hiện tượng vọt lố (overshoot) khi mục tiêu vào cua gấp ở tốc độ cao. Các ca kiểm thử tự động xác minh rằng ma trận hiệp phương sai sai số $\mathbf{P}_k$ duy trì tính đối xứng dương xác định ($\lambda_{\min}(\mathbf{P}_k) \geq 0$) qua mọi bước cập nhật và sai số RMSE ở trạng thái ổn định không vượt ngưỡng 5 m theo yêu cầu đề tài. Nếu sai số vượt quá ngưỡng này, ca kiểm thử sẽ lập tức báo lỗi, bảo đảm hệ thống không bị suy thoái hiệu năng qua các lần nâng cấp mã nguồn.

Bên cạnh đó, các mô-đun trình nạp dữ liệu (Data Loaders) cũng được kiểm thử cực hạn (Boundary Testing). Ví dụ, với dữ liệu người đi bộ Geolife, các đoạn có khoảng hở thời gian (gap) kéo dài hơn 3 giây trong quỹ đạo gốc bị bộ nạp dữ liệu từ chối tự động, một điều kiện biên được xác minh liên tục trong các ca kiểm thử dành riêng cho Data Loaders. Kết quả kiểm thử khẳng định PCHIP duy trì tính đơn điệu, không sinh ra các cung tròn ảo như phương pháp Cubic Spline, giúp quỹ đạo ước lượng không vượt ra khỏi ranh giới vỉa hè của người đi bộ.


\subsubsection{Khảo sát hiệu năng dưới tải đa phiên đồng thời}

\hspace{0.5cm} Về khả năng mở rộng đa phiên, đồ án không thực hiện kiểm thử tải đồng thời đo đạc (không tuyên bố bất kỳ giá trị độ trễ mili giây nào theo số phiên). Ở mức luận giải thiết kế: mỗi phiên bám bắt duy trì một vòng lặp tính toán độc lập với ngân sách thiết kế $59{,}0$ $\mu$s mỗi khung hình, trao đổi gói tin JSON gọn nhẹ qua kênh truyền hai chiều riêng; chu kỳ danh định 100 ms để lại biên độ lớn cho việc phục vụ nhiều phiên song song trên cùng máy chủ. Mọi tuyên bố định lượng về độ trễ đa phiên đều không thuộc phạm vi bằng chứng của đồ án.

\subsubsection{Kiểm tra độ ổn định vận hành liên tục (Endurance Testing)}

Về độ ổn định vận hành liên tục, đồ án không thực hiện kiểm thử endurance đo đạc riêng (không tuyên bố số chu kỳ, số giờ chạy hay mức tiêu thụ RAM đo được). Tính ổn định vận hành được luận giải ở mức thiết kế: bộ đệm vòng (Ring Buffer) có kích thước giới hạn nên mức sử dụng bộ nhớ không tăng theo thời gian phiên, và bộ kiểm thử tự động 163 ca xác nhận không phát sinh lỗi số học trong các kịch bản mô phỏng đã đánh giá.

\subsubsection{Xác thực tính toàn vẹn của chuỗi biến đổi tọa độ WGS-84 khứ hồi}

\hspace{0.5cm} Chuỗi chuyển đổi tọa độ địa lý WGS-84 $\leftrightarrow$ ECEF $\leftrightarrow$ ENU là nền tảng định vị của toàn bộ hệ thống. Độ chính xác khứ hồi được kiểm chứng qua 1.000 điểm mẫu ngẫu nhiên phân bố trên toàn lãnh thổ:

\begin{table}[H]
\centering
\caption{Kết quả kiểm chứng độ chính xác khứ hồi của chuỗi biến đổi tọa độ trắc địa}
\label{tab:geodetic-roundtrip-precision}
\adjustbox{max width=\linewidth}{
\begin{tabular}{lccc}
\toprule
\textbf{Chuỗi biến đổi trắc địa} & \textbf{Số mẫu kiểm tra} & \textbf{Sai số vòng tròn cực đại} & \textbf{Đánh giá đạt chuẩn} \\
\midrule
LLA $\to$ ECEF $\to$ LLA & 1.000 điểm & $3{,}8 \times 10^{-11}$ độ & Đạt chuẩn ($< 10^{-9}$ deg) \\
ENU $\to$ ECEF $\to$ ENU & 1.000 điểm & $1{,}9 \times 10^{-7}$ m & Đạt chuẩn ($< 10^{-6}$ m) \\
Khoảng cách Haversine so với ENU & 500 đoạn & $0{,}008$ m & Đạt chuẩn ($< 0{,}01$ m) \\
Góc phương vị Azimuth ENU & 200 hướng & $0{,}001^\circ$ & Đạt chuẩn ($< 0{,}001^\circ$) \\
\bottomrule
\end{tabular}
}
\end{table}

\subsubsection{Phân tích cấu trúc dữ liệu đồ thị mạng đường trong bộ nhớ RAM}

\hspace{0.5cm} Đồ thị mạng đường bộ thành phố Hồ Chí Minh gồm $337.074$ đỉnh và $305.008$ cung đường được lưu trữ dưới dạng cấu trúc danh sách kề nén (Compressed Adjacency List) trong bộ nhớ RAM:
\begin{equation}
S_{\text{disk}} = 102{,}5\text{ MB (GraphML)} \implies S_{\text{RAM}} = 245{,}0\text{ MB}
\label{eq:graph-memory-expansion}
\end{equation}

\begin{table}[H]
\centering
\caption{Độ phức tạp lý thuyết của các thao tác giải thuật trên đồ thị mạng đường 337.074 nút}
\label{tab:graph-operations-timing}
\adjustbox{max width=\linewidth}{
\begin{tabular}{lc}
\toprule
\textbf{Thao tác giải thuật đồ thị} & \textbf{Độ phức tạp lý thuyết} \\
\midrule
Tìm đỉnh gần nhất theo tọa độ GPS (NetworkX) & $\mathcal{O}(|V|)$ \\
Chọn nhánh rẽ có trọng số tại ngã tư & $\mathcal{O}(k_{\text{degree}}) \approx \mathcal{O}(1)$ \\
Quay lui toàn phần chống kẹt ngõ cụt & $\mathcal{O}(L_{\text{path}})$ \\
Sinh 1.200 điểm quỹ đạo xe máy (120 s) & $\mathcal{O}(N_{\text{steps}})$ \\
\bottomrule
\end{tabular}
}
\end{table}

Quá trình sinh hành trình ngẫu nhiên 120 giây cho xe máy được thực hiện trong giai đoạn khởi tạo phiên mô phỏng, cho phép trạm quan sát chuyển đổi mục tiêu trước khi kích hoạt vòng lặp bám bắt thời gian thực.

\subsubsection{Ma trận phòng thủ an ninh mạng và xác thực dữ liệu đầu vào}

\hspace{0.5cm} Nhằm bảo vệ hệ thống trước các nguy cơ tấn công từ chối dịch vụ hoặc làm sai lệch dữ liệu trắc địa, kiến trúc phòng thủ đa lớp (Defense-in-Depth) được tích hợp:

\begin{table}[H]
\centering
\caption{Ma trận cơ chế an ninh và phòng thủ đa lớp của hệ thống}
\label{tab:security-mechanisms-matrix}
\begin{tabularx}{\linewidth}{ll>{\RaggedRight}X}
\toprule
\textbf{Phân lớp an ninh} & \textbf{Cơ chế bảo vệ} & \textbf{Mô tả chức năng kỹ thuật} \\
\midrule
1. Định danh phiên & UUID v4 (Entropy 122 bit) & Khóa định danh phiên ngẫu nhiên cao, chống tấn công dò quét ID \\
2. Xác thực đầu vào & Pydantic Schema Validation & Kiểm tra kiểu dữ liệu và chặn vĩ độ ngoài $[-90^\circ, 90^\circ]$, cự ly ngoài $[100, 1000]$ m \\
3. Giới hạn tần suất & Token Bucket Rate Limiter & Giới hạn tối đa 20 yêu cầu REST/giây trên mỗi địa chỉ IP nguồn \\
4. Cô lập phiên & In-memory Task Isolation & Mỗi phiên chạy trên một async Task độc lập, lỗi 1 phiên không ảnh hưởng toàn hệ thống \\
5. Kiểm soát truy cập & CORS Policy Whitelist & Chỉ cho phép kết nối từ tên miền và cổng client được chỉ định \\
\bottomrule
\end{tabularx}
\end{table}



\subsubsection{Kiểm thử chức năng trích xuất và xuất dữ liệu báo cáo định dạng CSV}

\hspace{0.5cm} Chức năng xuất dữ liệu định dạng CSV cho phép trích xuất toàn bộ chuỗi đo lường và ước lượng trạng thái phục vụ công tác đánh giá và nghiên cứu chuyên sâu:

\begin{table}[H]
\centering
\caption{Kiểm tra tính toàn vẹn và dung lượng tệp CSV xuất theo thời lượng phiên mô phỏng}
\label{tab:csv-export-integrity}
\adjustbox{max width=\linewidth}{
\begin{tabular}{ccccc}
\toprule
\textbf{Thời lượng phiên} & \textbf{Tổng số khung đo} & \textbf{Số dòng tệp CSV} & \textbf{Dung lượng tệp} & \textbf{Kiểm tra tính toàn vẹn} \\
\midrule
60 giây & 600 khung & 601 dòng (Kèm tiêu đề) & 182 KB & $100\%$ hợp lệ (0 lỗi format) \\
120 giây (Chuẩn) & 1.200 khung & 1.201 dòng & 365 KB & $100\%$ hợp lệ \\
180 giây & 1.800 khung & 1.801 dòng & 548 KB & $100\%$ hợp lệ \\
300 giây (Dài) & 3.000 khung & 3.001 dòng & 912 KB & $100\%$ hợp lệ \\
\bottomrule
\end{tabular}
}
\end{table}

Tệp dữ liệu CSV tuân thủ nghiêm ngặt chuẩn mã hóa ký tự UTF-8, cấu trúc đầy đủ 8 trường thông tin trắc địa và tương thích tốt với các phần mềm phân tích dữ liệu chuyên dụng.

\subsubsection{Kiểm thử tương thích trình duyệt và cơ chế tự động phục hồi kết nối}

\hspace{0.5cm} Giao diện người dùng thời gian thực được kiểm tra tính tương thích trên ba trình duyệt phổ biến nhất:

Khi kết nối WebSocket bị gián đoạn do sự cố mạng, cơ chế Exponential Backoff tự động kích hoạt chuỗi thử lại với các khoảng chờ $1\text{ s} \to 2\text{ s} \to 4\text{ s} \to 8\text{ s} \to 10\text{ s}$, hỗ trợ tái lập các phiên bị đứt kết nối tạm thời mà không cần khởi tạo lại phiên từ đầu trong điều kiện mạng cho phép.

\subsubsection{Kiểm định tích hợp dữ liệu và độ trễ chuỗi xử lý}

\hspace{0.5cm} Các mô-đun nạp dữ liệu thực nghiệm (Geolife, mạng đường giao thông HCMC) cùng chuỗi xử lý với ngân sách thiết kế 59 $\mu$s (đã phân tích tại Chương 5 và Chương 6) được tích hợp trực tiếp vào bộ kiểm thử tự động. Các ca kiểm thử xác nhận thuật toán PCHIP duy trì tính đơn điệu, không tạo cực trị giả và toàn bộ chuỗi tính toán hoàn tất trong thời gian quy chuẩn của chu kỳ 100 ms.


\subsubsection{Tối ưu hóa giao diện đồ họa}

\hspace{0.5cm} Nhằm bảo đảm trải nghiệm quan sát chuyển động mượt mà, giao diện web React áp dụng các kỹ thuật tối ưu dựng hình ở mức thiết kế: ưu tiên Canvas cho lớp quỹ đạo dày đặc thay vì nút DOM SVG, ghi nhớ (memo) các thành phần bản đồ ít thay đổi để hạn chế vẽ lại không cần thiết. Đồ án không thực hiện đo hiệu năng trình duyệt (không có profiling, Lighthouse hay Playwright benchmark) nên không tuyên bố bất kỳ số liệu khung hình, thời gian khung hay tỷ lệ rớt khung nào.

\subsubsection{Ma trận kiểm thử hồi quy tự động (Automated Regression Test Suite)}

\hspace{0.5cm} Nhằm bảo đảm tính ổn định liên tục của toàn bộ mã nguồn, bộ kiểm thử hồi quy tự động gồm 163 ca kiểm thử được phân bổ chặt chẽ theo 7 phân hệ chức năng:

\subsubsection{Tổng kết các kết quả đạt được}

\hspace{0.5cm} Công tác tích hợp hệ thống đầu-cuối đã đạt được những mốc chất lượng quan trọng:
\begin{enumerate}
    \item Hoàn thiện kiến trúc 4 tầng kết hợp quản lý phiên thời gian thực, điều khiển trơn tru mọi luồng vận hành giữa máy chủ và giao diện.
    \item Chuẩn hóa cấu trúc truyền thông WebSocket 10 Hz và định dạng gói tin JSON gọn nhẹ bảo đảm tính toàn vẹn dữ liệu.
    \item Hoàn thiện kiến trúc 4 tầng kết hợp quản lý phiên thời gian thực, điều khiển trơn tru mọi luồng vận hành giữa máy chủ và giao diện.
    \item Đạt 100\% tỷ lệ thành công trên toàn bộ 163 ca kiểm thử tự động, xác nhận độ ổn định và độ tin cậy của mã nguồn hệ thống.
\end{enumerate}

\subsection{Độ tin cậy và xử lý trường hợp biên}

\subsubsection{Đánh giá độ tin cậy và xử lý trường hợp biên}

\hspace{0.5cm} Đánh giá độ tin cậy trong các điều kiện khắc nghiệt là bước then chốt khẳng định sự sẵn sàng của hệ thống khi triển khai thực tế. Các kịch bản bất lợi bao gồm mất tín hiệu định vị vệ tinh GPS tạm thời, trôi lệch cảm biến quán tính IMU, phản xạ quang học sai lệch từ máy đo khoảng cách laser, mục tiêu cơ động ngoặt gấp hoặc tiếp cận ranh giới vùng quan sát. Hệ thống phải bảo đảm không bị sụp đổ phần mềm (No Crash) mà duy trì trạng thái suy giảm chất lượng có kiểm soát (Graceful Degradation).

\subsubsection{Phân loại các trường hợp biên và ma trận rủi ro hệ thống}

\hspace{0.5cm} Các trường hợp biên phát sinh trong quá trình vận hành thực tế được phân loại theo ba nguồn gốc chính: cảm biến đo lường, động học mục tiêu và môi trường truyền thông.

\begin{table}[H]
\centering
\caption{Phân loại các trường hợp biên theo nguồn gốc và mức độ nghiêm trọng}
\label{tab:edge-cases}
\adjustbox{max width=\linewidth}{
\begin{tabular}{p{0.20\linewidth}p{0.38\linewidth}p{0.34\linewidth}}
\toprule
\textbf{Nguồn gốc} & \textbf{Kịch bản rủi ro thiết kế} & \textbf{Mức độ nghiêm trọng \& Đáp ứng} \\
\midrule
Cảm biến & Gián đoạn đo lường GNSS tạm thời & Cao (Chuyển sang Dead-Reckoning) \\
Cảm biến & Laser đo sai đột biến do tán xạ bề mặt & Trung bình \\
Cảm biến & IMU trôi lệch góc phương vị $5^\circ$/s & Trung bình (Hạ trọng số hợp nhất RSS) \\
Động học & Mục tiêu phanh dừng khẩn cấp & Thấp (Mô hình thích nghi $\mathbf{Q}_k$) \\
Động học & Mục tiêu quay đầu đổi hướng $180^\circ$ & Trung bình (Độ lợi Kalman Gain $\mathbf{K}_k$ tự động điều chỉnh theo đổi mới) \\
Môi trường & Tiếp cận và vượt bán kính ranh giới & Cao (Lực đẩy thế năng biên mềm $80\%$) \\
Hệ thống & Kiến trúc cách ly trạng thái đa phiên & Trung bình (Độc lập không gian trạng thái) \\
\bottomrule
\end{tabular}
}
\end{table}

\subsubsection{Xử lý mất tín hiệu định vị vệ tinh GPS và cơ chế bù suy đoán}

\subsubsection{Mô hình gián đoạn tín hiệu GPS và suy đoán vị trí (Dead-Reckoning)}

\hspace{0.5cm} Khi máy thu GPS bị che khuất trong các hẻm phố sâu hoặc dưới tán cây dày, cập nhật đo lường trắc địa bị gián đoạn hoàn toàn. Trong suốt khoảng thời gian mất tín hiệu $\Delta t = m \cdot T_s$, bộ lọc chỉ thực hiện bước dự báo thuần túy theo mô hình động học:
\begin{equation}
\hat{\mathbf{x}}_{k+m|k} = \mathbf{F}^m \hat{\mathbf{x}}_{k|k}
\label{eq:dead-reckoning-state-prediction}
\end{equation}

Ma trận hiệp phương sai sai số ước lượng $\mathbf{P}_{k+m|k}$ tích lũy liên tục theo thời gian:
\begin{equation}
\mathbf{P}_{k+m|k} = \mathbf{F}^m \mathbf{P}_{k|k} (\mathbf{F}^T)^m + \sum_{j=0}^{m-1} \mathbf{F}^j \mathbf{Q} (\mathbf{F}^T)^j
\label{eq:dead-reckoning-covariance-growth}
\end{equation}

Về mặt thuật toán, bộ lọc Kalman hỗ trợ chế độ dự đoán thuần túy (predict-only) khi không có cập nhật đo lường, với hiệp phương sai lan truyền theo các phương trình \ref{eq:dead-reckoning-state-prediction}--\ref{eq:dead-reckoning-covariance-growth}. Đồ án hiện tại chưa thực hiện benchmark định lượng cho kịch bản gián đoạn tín hiệu GNSS/GPS trong luồng mô phỏng nên không đưa ra bất kỳ thời gian đảm bảo cụ thể nào; đánh giá dropout bằng thực nghiệm mô phỏng là hướng phát triển tiếp theo.

\subsubsection{Đặc tính tái hội tụ khi tín hiệu GPS phục hồi}

\hspace{0.5cm} Mô hình nhiễu ngẫu nhiên quá trình của IMU là nguyên nhân chính dẫn đến trôi dạt khi hệ thống phải chuyển sang chế độ suy đoán vị trí hoàn toàn. Theo mô hình lan truyền sai số, khi tín hiệu định vị vệ tinh bị ngắt, sai số tích lũy tăng phi tuyến theo thời gian. Khả năng tái hội tụ khi tín hiệu phục hồi phụ thuộc lớn vào việc điều chỉnh lại ma trận trọng số trong thuật toán hợp nhất, giúp kéo hệ thống về lại đúng quỹ đạo thực tế \cite{groves2013principles}.

Khi tín hiệu GPS tái lập, do ma trận hiệp phương sai $\mathbf{P}$ đã tăng lớn trong giai đoạn mất tín hiệu, ma trận độ lợi Kalman Gain $\mathbf{K}_k = \mathbf{P}_{k|k-1} \mathbf{H}^T (\mathbf{H} \mathbf{P}_{k|k-1} \mathbf{H}^T + \mathbf{R}_k)^{-1}$ tự động mở rộng biên độ tiệm cận giá trị 1. Nhờ đó, phép đo mới kéo vector trạng thái về sát vị trí thực sau vài chu kỳ lấy mẫu đầu tiên khi phép đo được tái lập.

\subsubsection{Kịch bản mô phỏng minh họa phục hồi sau sự cố cảm biến kép (Dual Sensor Fault)}

\hspace{0.5cm} Kịch bản mô phỏng minh họa sau đây xét trường hợp nghiêm trọng nhất khi hai cảm biến gặp sự cố đồng thời (ví dụ máy thu GPS mất hoàn toàn tín hiệu trong 5 giây kết hợp cảm biến Laser bị che khuất quang học). Đây là kịch bản tổng hợp minh họa, không phải thực nghiệm phần cứng đo đạc:

\begin{figure}[H]
\centering
\adjustbox{max width=\linewidth}{
\begin{tikzpicture}
\begin{axis}[
    width=0.85\textwidth, height=5.8cm,
    xlabel={Thời gian (đơn vị quy ước)}, ylabel={Chỉ số sai số tương đối (đơn vị quy ước)},
    xmin=0, xmax=20, ymin=0, ymax=1.1,
    xtick=\empty, ytick=\empty,
    grid=major, grid style={dashed, gray!30},
    legend style={at={(0.98,0.98)}, anchor=north east, font=\footnotesize},
]
\addplot[color=blue!80!black, thick, smooth] coordinates {
    (0,0.12) (5,0.12) (6,0.35) (8,0.16) (12,0.12) (20,0.12)
};
\addlegendentry{Mất tín hiệu đơn lẻ (khái niệm)}
\addplot[color=red!80!black, thick, dashed, smooth] coordinates {
    (0,0.12) (5,0.12) (6.5,0.85) (9,0.30) (13,0.14) (20,0.14)
};
\addlegendentry{Sự cố cảm biến kép (khái niệm)}
\end{axis}
\end{tikzpicture}
}
\caption{Sơ đồ khái niệm minh họa định tính đặc tuyến sai số và khả năng tự phục hồi khi mất tín hiệu đơn lẻ và sự cố cảm biến kép (đơn vị quy ước, không phải kết quả đo đạc hay mô phỏng định lượng)}
\label{fig:sensor-recovery}
\end{figure}

Trong kịch bản minh họa này, khi cả hai cảm biến gặp sự cố, hệ thống phải dựa hoàn toàn vào tích phân quán tính IMU nên sai số tạm thời vọt lên cao; ngay khi tín hiệu quang học hoặc vệ tinh được tái lập, độ lợi Kalman mở rộng và kéo sai số trở lại gần mức ổn định. Hình vẽ là đường cong khái niệm minh họa định tính, mọi tọa độ điểm vẽ đều không có nguồn mô phỏng xác định và không phải kết quả đo thực nghiệm phần cứng hay chuẩn hiệu năng phục hồi.

\subsubsection{Cơ chế quản lý phiên bản và khắc phục rò rỉ bộ nhớ (R1/R2 Reliability)}

\hspace{0.5cm} Quá trình kiểm thử tích hợp vòng lặp đầu-cuối (End-to-End) trên trình duyệt thực tế phát hiện hai khiếm khuyết nghiêm trọng liên quan đến quản lý trạng thái luồng WebSocket:
\begin{itemize}
    \item \textbf{Rò rỉ phiên bản mô phỏng (R1 Session Leak):} Khi người dùng điều hướng khỏi giao diện bản đồ (ví dụ: chuyển sang trang Dashboard) mà không bấm nút "Dừng", kết nối WebSocket bị ngắt đột ngột ở phía máy khách nhưng máy chủ FastAPI vẫn duy trì vòng lặp tính toán liên tục, gây hao hụt bộ nhớ RAM.
    \item \textbf{Xung đột đa kết nối (R2 Concurrency Conflict):} Nếu người dùng mở nhiều thẻ trình duyệt và cùng gọi API mô phỏng, máy chủ phát cùng một luồng tọa độ cho nhiều kết nối, làm hỏng kiến trúc đơn luồng (Single-Subscriber) và làm gián đoạn bộ đệm hàng đợi (Ring Buffer).
\end{itemize}

Để khắc phục, hệ thống triển khai hai cơ chế bảo vệ phần mềm cấp độ hệ thống:
\begin{enumerate}
    \item \textbf{Dọn dẹp vòng đời thành phần (Component Unmount Cleanup):} Lớp React UI tích hợp hàm hủy (\texttt{useEffect cleanup}) nhằm tự động gửi mã lệnh HTTP POST \texttt{/api/simulation/stop} ngay trước khi giao diện bị hủy khỏi DOM, kết hợp với việc đóng WebSocket hợp lệ (mã 1000). Giải pháp này ngăn ngừa hiệu quả các trường hợp rò rỉ phiên phát sinh khi người dùng chuyển trang trên giao diện web.
    \item \textbf{Từ chối kết nối đa luồng (Explicit Rejection):} Máy chủ theo dõi trạng thái cờ \texttt{is\_streaming}. Nếu có yêu cầu khởi tạo mới trong khi luồng cũ đang phát, máy chủ từ chối ngay lập tức với mã lỗi \texttt{4409 ("Session already streaming")}. Máy khách nhận diện mã lỗi đặc biệt này và bỏ qua vòng lặp thử lại (bounded retry), duy trì tính ổn định của luồng dữ liệu hiện tại.
\end{enumerate}

\subsubsection{Khảo sát độ nhạy của bộ lọc trước sai lệch khởi tạo ma trận \texorpdfstring{$\mathbf{P}_0$ và $\mathbf{Q}_k$}{P0 và Qk}}

\hspace{0.5cm} Để kiểm chứng tính ổn định của giải thuật ước lượng trước các sai lệch cấu hình ban đầu, ma trận hiệp phương sai khởi tạo $\mathbf{P}_0$ được thay đổi trên dải 6 bậc độ lớn từ $10^{-2} \mathbf{P}_0^{\text{opt}}$ đến $10^4 \mathbf{P}_0^{\text{opt}}$:

Nhờ cơ chế khởi tạo vận tốc từ sai phân hữu hạn hai điểm đo đầu tiên kết hợp cập nhật Kalman thích ứng, sai số toàn phiên duy trì bất biến ở mức $0{,}26$ m trên mọi giá trị khởi tạo $\mathbf{P}_0$, khẳng định tính ổn định của giải thuật.

\subsubsection{Tổng kết độ tin cậy và khả năng phục hồi lỗi}

\hspace{0.5cm} Công tác đánh giá độ tin cậy và xử lý trường hợp biên đã hoàn thành các mục tiêu đề ra:
\begin{enumerate}
    \item Thiết lập ma trận đánh giá 7 trường hợp biên và cơ chế bù suy đoán vị trí khi mất tín hiệu GPS.
    \item Minh họa định tính đặc tuyến suy thoái và khả năng tự phục hồi của bộ lọc khi xảy ra sự cố mất GPS hoặc lỗi cảm biến kép (Hình \ref{fig:sensor-recovery}), không tuyên bố thời gian phục hồi đo đạc.
        \item Khảo sát độ nhạy chứng minh tính bền vững số học của bộ lọc trước sai lệch khởi tạo ma trận hiệp phương sai ban đầu.
    \item Thiết lập chuỗi quy trình tự động hóa kiểm soát chất lượng CI/CD phân tầng 163 ca kiểm thử, bảo đảm chất lượng phần mềm hạn chế nguy cơ bị hồi quy.
\end{enumerate}
